\subsection{The graph identity and its coercive quadratic form}
\label{sec:fl-alg}

We supply the pointwise calculations used in the flat deformation.
All derivatives, norms, traces, contractions, and divergences in this
subsection refer to \(g\), unless a different metric is displayed.
A one-form in a divergence is identified with its \(g\)-dual vector.
For symmetric covariant tensors define
\[
 [A,B]=\langle A,B\rangle-(\tr A)(\tr B),
 \qquad p_A=A-(\tr A)g.
\]
In particular, the flat constraints have the normalization
\begin{equation}
 R_g=16\pi\mu+[K,K],
 \qquad \operatorname{div}p_K=8\pi J.
 \label{eq:fl-alg-constraints}
\end{equation}
There is no cosmological term in these identities.

Write \(p=(\log l)'(t)\), so that \(\dd l=pl\,\dd t\), and retain the
graph and trace quantities defined above.
The identities in this subsection require only \(l>0\), \(0\le p\le N\),
and the trace equation \(\tr_{A_\chi}S=F\).
At a point where \(\dd f\ne0\), use an orthonormal frame with first
vector \(e=\nabla f/|\dd f|\), and decompose
\[
 S=\begin{pmatrix}q_1&M^{\mathsf T}\\M&P\end{pmatrix},
 \qquad \dd t=x e^\flat+y,\qquad y(e)=0.
\]
Here \(P\) is a symmetric tensor on the two-dimensional plane
\(e^\perp\), and
\begin{equation}
 v=a^2=1-d,\qquad
 \chi=\frac{4d}{4+pv},\qquad
 P_0:=\tr P=F-\chi q_1.
 \label{eq:fl-alg-blocks}
\end{equation}
We use \(P^{\mathrm{tf}}=P-\frac12(\tr P)g|_{e^\perp}\).
The symbol \(|\dd t|_{A_d}^2\) means
\(\langle A_d\nabla t,\nabla t\rangle=dx^2+|y|^2\).

\begin{lemma}[Differentiation and flux identities]
\label{lem:fl-alg-flux}
For \(U=l\sqrt D\), one has
\begin{align}
 \tfrac12\nabla\log D
   &=H^f(w,\cdot)+pv\,\dd t,
   \label{eq:fl-alg-D-derivative}\\
 \nabla_iw_j
   &=(H^f A_d)_{ij}+pd\,(\dd t)_i w_j,
   \label{eq:fl-alg-w-derivative}\\
 \nabla\log u
   &=H^f(w,\cdot)+(2+p+pv)\dd t.
   \label{eq:fl-alg-u-derivative}
\end{align}
Set
\[
 k=-H^f-p(\dd t\otimes w^\flat+w^\flat\otimes\dd t),
 \qquad Q=K-k.
\]
Then the flux \(V=uA_\chi(4\nabla Z+K(w,\cdot))\) satisfies
\begin{align}
 \frac{V}{u}
    &=p_Q(w,\cdot)+4A_d\nabla t+wF,
    \label{eq:fl-alg-flux-identity}\\
 \frac{\operatorname{div}(uw)}u
    &=F+(1-\chi)q_1-\tr K+2(p+1)ax.
    \label{eq:fl-alg-uw-divergence}
\end{align}
Every expression asserted to be an identity extends smoothly across
\(\dd f=0\).
\end{lemma}

\begin{proof}
Differentiate \(D=1+l^2|\dd f|^2\), use \(\dd l=pl\,\dd t\), and
substitute \(H^f=l\Hess f/\sqrt D\) and
\(w=l\nabla f/\sqrt D\). This gives
\eqref{eq:fl-alg-D-derivative}. Differentiating the last formula for
\(w\), and using \(1-v=d\), gives
\[
 \nabla_iw_j
 =H^f_{ij}-H^f(w,\cdot)_i w_j+pd\,(\dd t)_i w_j.
\]
Since \(A_d=I-w\otimes w^\flat\), this is
\eqref{eq:fl-alg-w-derivative}.
Equation \eqref{eq:fl-alg-u-derivative} follows from
\(u=e^{2t}l\sqrt D\).
Taking the trace of \eqref{eq:fl-alg-w-derivative} and adding
\(\langle\nabla\log u,w\rangle\) gives
\[
 \frac{\operatorname{div}(uw)}u
 =\tr H^f+2(p+1)ax.
\]
Now \(\tr H^f=F+(1-\chi)q_1-\tr K\), proving
\eqref{eq:fl-alg-uw-divergence}.

From \(Z=t+\frac18\log D\), Equation
\eqref{eq:fl-alg-D-derivative} yields
\[
 4\nabla Z+K(w,\cdot)=(4+pv)\nabla t+S(w,\cdot).
\]
Its image under \(A_\chi\) has axial and transverse components
\[
 4dx+\chi aq_1,\qquad (4+pv)y+aM,
\]
because \(\chi(4+pv)=4d\).
On the other hand,
\[
 Q_{ee}=q_1+2pa x,\qquad Q_{e\perp}=M+pa y,
 \qquad Q_\perp=P.
\]
The same two components of
\(p_Q(w,\cdot)+4A_d\nabla t+wF\) are
\(-aP_0+4dx+aF\) and \(a(M+pa y)+4y\).
Equation \eqref{eq:fl-alg-blocks} identifies these with the preceding
display, proving \eqref{eq:fl-alg-flux-identity}.
Smoothness follows either from the original definitions or from
\begin{equation}
 A_\chi
 =I-\frac{p+4}{4+p|w|^2}\,w\otimes w^\flat.
 \label{eq:fl-alg-smooth-matrix}
\end{equation}
In particular no choice of \(e\) is needed at \(w=0\).
\end{proof}

\begin{proposition}[Scalar-curvature identity]
\label{prop:fl-alg-scalar}
Let
\[
 \bar g=g+l^2\dd f^2,\qquad \hat g=e^{4t}\bar g.
\]
For \(\mathcal T\) in \eqref{eq:fl-quadratic-form}, the exact identity
is
\begin{equation}
 \frac12 e^{4t}R_{\hat g}
 =8\pi\bigl(\mu+J(w)\bigr)+\mathcal T+w(F)
       -F\tr K-\frac{\operatorname{div}V}{u}.
 \label{eq:fl-alg-identity}
\end{equation}
\end{proposition}

\begin{proof}
\emph{The stationary calculation.}
On a product with coordinate \(s_0\), introduce the stationary Lorentz
metric
\[
 \mathbf g=-l^2(\dd s_0-\dd f)^2+\bar g
           =-l^2\dd s_0^2+2l^2\dd s_0\dd f+g.
\]
The \(s_0\)-slices have induced metric \(g\), shift
\(\beta=l^2\nabla f=Uw\), lapse \(U=l\sqrt D\), and unit normal
\(n=U^{-1}(\partial_{s_0}-\beta)\). We choose the sign
\(k(X,Y)=\mathbf g(\mathbf\nabla_Xn,Y)\).
Since the slice metric is independent of \(s_0\),
\[
 k=-\frac{\operatorname{sym}\nabla\beta}{U}
   =-H^f-p(\dd t\otimes w^\flat+w^\flat\otimes\dd t).
\]
Here \(\operatorname{sym}T=(T+T^{\mathsf T})/2\).

For completeness, put \(\vartheta=\tr k\). The contracted Gauss
equation and the normal expansion equation, with this sign convention,
read
\[
 R_{\mathbf g}=R_g-2\Ric_{\mathbf g}(n,n)-|k|^2+\vartheta^2,
 \qquad
 \Ric_{\mathbf g}(n,n)=-n(\vartheta)-|k|^2+U^{-1}\Delta_gU.
\]
The lapse term in the second formula follows from the acceleration
\(\mathbf\nabla_n n=\nabla\log U\); its divergence contribution is
\(\Delta_g\log U+|\nabla\log U|^2=U^{-1}\Delta_gU\).
Stationarity gives \(Un(\vartheta)=-\beta(\vartheta)\), while
\(\operatorname{div}\beta=-U\vartheta\).
Combining these equalities gives
\begin{equation}
 UR_{\mathbf g}
 =U(R_g+[k,k])
   -2\operatorname{div}\bigl(\nabla U+\vartheta\beta\bigr).
 \label{eq:fl-alg-stationary-scalar}
\end{equation}

In the coordinate \(s=s_0-f\), the same metric is
\(-l^2\dd s^2+\bar g\). Direct contraction of its warped-product
curvature gives
\(R_{\mathbf g}=R_{\bar g}-2l^{-1}\bar\Delta l\).
The determinant and inverse metric of \(\bar g\) give, for a scalar
\(\varphi\),
\[
 \bar\Delta\varphi
 =D^{-1/2}\operatorname{div}(\sqrt D\,A_d\nabla\varphi).
\]
Moreover, Equations \eqref{eq:fl-alg-D-derivative} and
\eqref{eq:fl-alg-w-derivative} imply
\[
 \nabla U-\frac Ul A_d\nabla l=-k(\beta,\cdot).
\]
For example, after division by \(U\) its left side is
\[
 H^f(w,\cdot)+p\{v\,\dd t+w^\flat\langle w,\nabla t\rangle\}
 =-k(w,\cdot).
\]
Substituting into \eqref{eq:fl-alg-stationary-scalar} therefore yields
\begin{equation}
 R_{\bar g}=R_g+[k,k]+\frac2U\operatorname{div}(p_k\beta).
 \label{eq:fl-alg-graph-scalar}
\end{equation}

\emph{The constraint substitution.}
The symmetric part of \(\nabla\beta\) is \(-Uk\), so
\eqref{eq:fl-alg-constraints} gives
\[
 \operatorname{div}(p_K\beta)
 =8\pi UJ(w)-U[K,k].
\]
Use \(Q=K-k\), \(p_Q=p_K-p_k\), and
\([Q,Q]=[K,K]+[k,k]-2[K,k]\) in
\eqref{eq:fl-alg-graph-scalar}. The result is
\begin{equation}
 R_{\bar g}
 =16\pi(\mu+J(w))+[Q,Q]
      -\frac2U\operatorname{div}(Up_Q(w,\cdot)).
 \label{eq:fl-alg-constraint-graph}
\end{equation}

\emph{The conformal calculation.}
In dimension three,
\[
 \frac12 e^{4t}R_{\hat g}
 =\frac12R_{\bar g}-4\bar\Delta t-4|\dd t|_{A_d}^2.
\]
Replacing the weight \(U\) in the last flux of
\eqref{eq:fl-alg-constraint-graph} by \(u=e^{2t}U\) adds
\(2p_Q(w,\nabla t)\).
Writing \(\bar\Delta t\) in the preceding divergence form, and using
\(\log(u/\sqrt D)=2t+\log l\), gives
\[
 -4\bar\Delta t-4|\dd t|_{A_d}^2
 =-\frac1u\operatorname{div}(4uA_d\nabla t)
        +4(p+1)|\dd t|_{A_d}^2.
\]
By \eqref{eq:fl-alg-flux-identity}, the two fluxes obtained so far
sum to \(V-uwF\). Its divergence contributes
\[
 \frac1u\operatorname{div}(uwF)
 =w(F)+F\frac{\operatorname{div}(uw)}u.
\]
It follows that the quadratic term remaining after the displayed
\(-F\tr K\) in \eqref{eq:fl-alg-identity} is
\begin{equation}
 \frac12[Q,Q]+2p_Q(w,\nabla t)+4(p+1)|\dd t|_{A_d}^2
       +F\{F+(1-\chi)q_1+2(p+1)ax\}.
 \label{eq:fl-alg-invariant-quadratic}
\end{equation}
The next calculation identifies it with \(\mathcal T\), completing
the proof.
\end{proof}

\begin{lemma}[Square completion and uniformity]
\label{lem:fl-alg-coercivity}
The expression \eqref{eq:fl-alg-invariant-quadratic} equals
\(\mathcal T\) in \eqref{eq:fl-quadratic-form}. It is smooth in the
actual variables at \(\dd f=0\). For an absolute constant \(C\), if
\(0\le p\le N\), then
\begin{equation}
 F^2+|P|^2+|M|^2+dq_1^2+|\dd t|_{A_d}^2
       \le C(N+4)^2\mathcal T.
 \label{eq:fl-alg-general-coercivity}
\end{equation}
At a point with \(\dd t=0\), put
\(\mathcal E=|P^{\mathrm{tf}}|^2+|M|^2+F^2+\chi q_1^2\).
There is the stronger bound
\begin{equation}
 \frac12\mathcal E\le\mathcal T\le\mathcal E.
 \label{eq:fl-alg-floor-comparison}
\end{equation}
These constants are independent of \(d\), \(p\), \(N\), \(l\), and
the size of \(\dd f\).
\end{lemma}

\begin{proof}
Using \(P_0=F-\chi q_1\), direct contraction gives
\begin{align*}
 \frac12[Q,Q]
 &=\frac12|P^{\mathrm{tf}}|^2-\frac14P_0^2
   +|M+pa y|^2-P_0(q_1+2pa x),\\
 2p_Q(w,\nabla t)&=-2aP_0x+2aM\cdot y+2pv|y|^2.
\end{align*}
Consequently \eqref{eq:fl-alg-invariant-quadratic} is
\begin{align}
 \mathcal T={}&\frac12|P^{\mathrm{tf}}|^2
     +|M+(p+1)a y|^2+[4(p+1)-v]|y|^2
       +4(p+1)dx^2+2(p+1)a\chi q_1x \notag\\
 &\quad+\frac34F^2-\frac12\chi Fq_1
                +(\chi-\chi^2/4)q_1^2 .
 \label{eq:fl-alg-expanded-quadratic}
\end{align}
Completion of its axial and trace squares is exact because
\[
 \chi-\frac{\chi^2}{4}
 =\frac{(p+1)\chi^2v}{4d}+\frac{\chi^2}{12}
                  +\frac{4d(9-d)}{3(4+pv)^2}.
\]
This proves the asserted formula for \(\mathcal T\).

For smoothness, rewrite
\((1-\chi)q_1=\tr S-F\) in
\eqref{eq:fl-alg-invariant-quadratic}, and
\(ax=\langle w,\nabla t\rangle\).
Every term is then an invariant smooth expression in \(S,w,\dd t,p\),
with \(F=\tr_{A_\chi}S\) and \(A_\chi\) given by
\eqref{eq:fl-alg-smooth-matrix}.
At \(w=0\) its value is
\[
 \mathcal T=\frac12|S|^2+\frac12(\tr S)^2
                               +4(p+1)|\dd t|^2.
\]
This also defines the direction-independent value at a critical point
of \(f\).

To prove \eqref{eq:fl-alg-general-coercivity}, denote the translated
variables in the axial, trace, and mixed squares by
\[
 z=x+\frac{\chi a q_1}{4d},\qquad
 r_F=F-\frac{\chi q_1}{3},\qquad
 m_F=M+(p+1)a y.
\]
The square formula bounds \(|P^{\mathrm{tf}}|^2,|m_F|^2\),
\((p+1)|y|^2\), \((p+1)dz^2\), and \(r_F^2\) by absolute multiples
of \(\mathcal T\). Its last coefficient satisfies
\[
 \frac{4d(9-d)}{3(4+pv)^2}
                    \ge \frac{32d}{3(N+4)^2},
\]
so \(dq_1^2\le C(N+4)^2\mathcal T\).
Since \(\chi\le d\le1\), one has
\[
 dx^2\le 2dz^2+\frac{\chi^2v}{8d}q_1^2
                    \le2dz^2+\frac18dq_1^2,
 \qquad
 F^2\le2r_F^2+\frac29\chi^2q_1^2.
\]
These estimates control \(x,F\), and
\(|P|^2=|P^{\mathrm{tf}}|^2+\frac12(F-\chi q_1)^2\).
Finally,
\[
 |M|^2\le2|m_F|^2+2(p+1)^2v|y|^2
                               \le C(N+1)\mathcal T.
\]
They prove \eqref{eq:fl-alg-general-coercivity}.

If \(\dd t=0\), Equation \eqref{eq:fl-alg-expanded-quadratic} becomes
\[
 \mathcal T=\frac12|P^{\mathrm{tf}}|^2+|M|^2
       +\frac34F^2-\frac12\chi Fq_1
                      +\chi(1-\chi/4)q_1^2.
\]
Using
\(\frac12|\chi Fq_1|\le\frac14F^2+\frac14\chi^2q_1^2\)
and \(0<\chi\le1\) proves both bounds in
\eqref{eq:fl-alg-floor-comparison}.
\end{proof}

\subsection{Pointwise estimates at an interior floor contact}
\label{sec:fl-alg-floor}

A floor contact below refers to an interior local minimum of \(t\).
Thus \(\dd t=0\), \(\Hess t\ge0\), and
\(\Hess l/l=p\Hess t\ge0\).
The estimates hold at any such minimum, not only at the particular
height chosen in the continuation.

\begin{lemma}[Gauss Inequality at a floor contact]
\label{lem:fl-alg-floor-gauss}
For a symmetric tensor \(A\), define
\[
 \mathcal S(A)
 =\frac14(\tr_\perp A)^2-\frac12|A_\perp^{\mathrm{tf}}|^2
       +dA_{ee}\tr_\perp A-d|A_{e\perp}|^2.
\]
At an interior floor contact,
\begin{equation}
 \frac12e^{4t}R_{\hat g}
 \le \frac12R_g-v\Ric_g(e,e)+\mathcal S(H^f).
 \label{eq:fl-alg-floor-gauss}
\end{equation}
The expression \(v\Ric_g(e,e)\) is interpreted as \(\Ric_g(w,w)\);
all quantities have continuous, direction-independent values at
\(w=0\).
\end{lemma}

\begin{proof}
Consider the graph of \(f\) in the Riemannian warped product
\(g+l^2\dd s^2\). Its induced metric is \(\bar g\), its unit normal
has horizontal component \(-w\), and its vertical component has
length \(\sqrt d\).
Write \(T=l^{-1}\partial_s\) for the unit vertical vector.
The ambient curvature formulas are
\begin{align*}
 R_{\mathrm{amb}}&=R_g-2l^{-1}\Delta_gl,\\
 \Ric_{\mathrm{amb}}(X,Y)&=\Ric_g(X,Y)-l^{-1}\Hess l(X,Y),\\
 \Ric_{\mathrm{amb}}(T,T)&=-l^{-1}\Delta_gl,
\end{align*}
with zero mixed Ricci components.
Hence half the ambient term
\(R_{\mathrm{amb}}-2\Ric_{\mathrm{amb}}(n,n)\) in the graph Gauss
equation is
\[
 \frac12R_g-v\Ric_g(e,e)
                         -\frac vl\tr_\perp\Hess l.
\]
At the contact \(\dd l=0\); the covariant components of the graph
second form are therefore \(H^f\), up to the harmless common sign
coming from the choice of graph normal.
Tracing with the graph inverse metric \(A_d\) gives
\[
 \frac12\bigl((\tr_{A_d}A)^2-|A|_{A_d}^2\bigr)=\mathcal S(A).
\]
Indeed the axial square cancels, the mixed contribution is
\(-d|A_{e\perp}|^2\), and the transverse contribution is
\(\frac14(\tr_\perp A)^2-\frac12|A_\perp^{\mathrm{tf}}|^2\).
Finally, at \(\dd t=0\) the conformal scalar formula contributes
\(-4\tr(A_d\Hess t)\). We have thus proved the exact formula
\[
 \frac12e^{4t}R_{\hat g}
 =\frac12R_g-v\Ric_g(e,e)+\mathcal S(H^f)
       -\frac vl\tr_\perp\Hess l-4\tr(A_d\Hess t).
\]
Both final terms are nonpositive, proving
\eqref{eq:fl-alg-floor-gauss}.
\end{proof}

\begin{lemma}[Uniform floor coercivity]
\label{lem:fl-alg-floor-coercivity}
At a point with \(\dd t=0\), for an absolute constant \(C\),
\begin{equation}
 \mathcal T-\mathcal S(H^f)
 \ge \frac1{16}\mathcal T
        -C(1+N)^2\bigl(|K|^2+vF^2\bigr).
 \label{eq:fl-alg-floor-coercivity}
\end{equation}
In particular the coefficient \(1/16\) is independent of \(N\) and
of the graph gradient.
\end{lemma}

\begin{proof}
Subtracting \(\mathcal S(S)\) from
\eqref{eq:fl-alg-expanded-quadratic}, with \(\dd t=0\), gives
\begin{align}
 \mathcal T-\mathcal S(S)
  ={}&|P^{\mathrm{tf}}|^2+(1+d)|M|^2
       +\chi(1+d-\chi/2)q_1^2-dFq_1+\frac12F^2.
 \label{eq:fl-alg-floor-difference}
\end{align}
The identities
\[
 d\ge\chi,\qquad d-\chi=\frac{pv}{4}\chi,\qquad
 \frac{d^2}{\chi}=\frac{d(4+pv)}4\le\frac{N+4}{4}
\]
make its parameter dependence explicit.
For the part with cross coefficient \(-\chi\), Young's Inequality gives
\[
 \frac12F^2-\chi Fq_1+\chi(1+\chi/2)q_1^2
       \ge\frac18F^2+\frac56\chi q_1^2.
\]
The excess cross term satisfies
\[
 |(d-\chi)Fq_1|
 \le\frac13\chi q_1^2
           +\frac34\frac{(d-\chi)^2}{\chi}F^2
 \le\frac13\chi q_1^2+\frac3{64}p^2vF^2.
\]
It follows from \eqref{eq:fl-alg-floor-comparison} that
\begin{equation}
 \mathcal T-\mathcal S(S)
       \ge\frac18\mathcal T-\frac3{64}p^2vF^2.
 \label{eq:fl-alg-unshifted-floor}
\end{equation}

Because \(H^f=S-K\), expansion of the quadratic polynomial
\(\mathcal S\) gives
\[
 |\mathcal S(S-K)-\mathcal S(S)|
 \le C|K|\bigl(|P^{\mathrm{tf}}|+|M|+|F|+d|q_1|\bigr)+C|K|^2.
\]
Here we used \(\tr P=F-\chi q_1\) and \(\chi\le d\).
The estimate for \(d^2/\chi\) above and
\eqref{eq:fl-alg-floor-comparison} bound this by
\[
 \frac1{16}\mathcal T+C(1+N)|K|^2.
\]
Combining with \eqref{eq:fl-alg-unshifted-floor} proves the claim.
\end{proof}

\subsection{Differentiating the trace and the weighted tensor drift}
\label{sec:fl-alg-derivatives}

The next lemma records precisely the hypotheses on the continuation
right sides that enter the floor argument. They concern values and
first derivatives only. Let \(E\) be a fixed compact set containing
the support of \(\tr K\) and of all collar corrections. Fix
\(0<\delta<1\) and \(3/4<\gamma<1\), and positive smooth weights
\(\rho_0,\rho_1\) with
\[
 \rho_0\asymp r^{-3-\delta},\qquad
 \rho_1\asymp r^{-2\gamma}
\]
on the Euclidean end. Suppose
\[
 |K|\le Cr^{-2},\qquad |\nabla K|\le Cr^{-3}
\]
there. All constants defining \(E\), the weights, and these bounds
are fixed before \(N\).

For clarity denote the trace right side as a function of its inputs
by \(\mathfrak F(x,w,h,p)\), and its value on the current fields by
\(F\). The allowed forms are
\begin{align}
 \mathfrak F_1&=h+C(x), \notag\\
 \mathfrak F_{2,\theta}
    &=h+C(x)+\theta D_0(x,w,h), \qquad 0\le\theta\le1,\notag\\
 \mathfrak F_{3,\zeta}
    &=h+(1-\zeta)\{C(x)+D_0(x,w,h+\zeta b(x))\}
       +\zeta\{\tr_{A_\chi}K+D_1(x,w)\},
       \qquad 0\le\zeta\le1 .
 \label{eq:fl-alg-stage-right-sides}
\end{align}
Here \(C,b\) are fixed smooth functions, and
\(D_0,D_1\) are smooth bundle functions, supported over \(E\), that
extend smoothly to \(|w|\le1\).
Assume \(\partial_hD_0\ge0\).
On the bounded height interval under consideration, their values and
first derivatives in the base, fiber, and height variables are
bounded independently of \(N\).
Only \(A_\chi\), through \(p\), has \(N\)-dependent coefficients.
Base derivatives of bundle functions mean covariant derivatives with
their fiber arguments parallel transported.

\begin{lemma}[Floor derivative estimates]
\label{lem:fl-alg-floor-derivatives}
Assume the preceding hypotheses, \(h=\tau f\), and
\(|h|\le C_0\). At any point where \(\dd t=0\), for each fixed
\(\eta>0\) there are constants \(C_\eta\) and an integer \(m\),
independent of \(N\), the homotopy parameters, and the graph gradient,
such that
\begin{align}
 w(F)&\ge \tau a|\dd f|-\eta\mathcal T
              -C_\eta(1+N)^m(\rho_0+v\rho_1),
       \label{eq:fl-alg-trace-derivative}\\
 \left|\frac{\operatorname{div}(uA_\chi K(w,\cdot))}{u}\right|
      &\le\eta\mathcal T+C_\eta(\rho_0+v\rho_1).
       \label{eq:fl-alg-drift-derivative}
\end{align}
If in addition \(|h|\le Cr^{-\gamma}\) on the end, then
\begin{equation}
 |K|^2+vF^2\le C(\rho_0+v\rho_1).
 \label{eq:fl-alg-weighted-remainder}
\end{equation}
The constants may depend on the fixed data, weights, compact
corrections, height bound, and \(\eta\); they are uniform over the
three stages in \eqref{eq:fl-alg-stage-right-sides}.
\end{lemma}

\begin{proof}
\emph{The error weights.}
Since \(\delta+2\gamma<3\), the end decay implies
\[
 |\nabla K|\le C\sqrt{\rho_0\rho_1},\qquad
 |K|^2\le C\rho_0.
\]
Young's Inequality therefore gives
\begin{equation}
 a|\nabla K|\le C(\rho_0+v\rho_1).
 \label{eq:fl-alg-gradient-weight}
\end{equation}
Any bounded error supported in \(E\) is bounded by \(C\rho_0\),
because \(\rho_0\) has a positive minimum there.

\emph{The exact tensor-drift derivative.}
Put \(H=H^f\), \(k_0=K_{ee}\), \(m_0=K_{e\perp}\), and
\(P_H=H|_{e^\perp}\). At the point in question,
\[
 \nabla p=0,\qquad \nabla\log u=H(w,\cdot),
 \qquad \nabla w=H A_d.
\]
For \(c=(p+4)/(4+pv)\), Equation
\eqref{eq:fl-alg-smooth-matrix} reads
\[
 A_\chi K(w,\cdot)=K(w,\cdot)-cK(w,w)w.
\]
Its exact differentiated form is
\begin{align}
 \frac{\operatorname{div}(uA_\chi K(w,\cdot))}{u}
 ={}&a\{(\operatorname{div}K)(e)
                     -(1-\chi)(\nabla_eK)(e,e)\}\notag\\
 &+\chi(2\chi-1)H_{ee}k_0
      +2\chi\langle H_{e\perp},m_0\rangle
      +P_H:K_\perp-(1-\chi)(\tr P_H)k_0.
 \label{eq:fl-alg-exact-drift}
\end{align}
To see the cancellation explicitly, use a normal orthonormal frame
at the point. The entries of \(\nabla w=H A_d\) are
\[
 \nabla_e w_e=dH_{ee},\quad
 \nabla_e w_A=H_{eA},\quad
 \nabla_A w_e=dH_{Ae},\quad
 \nabla_A w_B=H_{AB}.
\]
The coefficient of \(H_{ee}k_0\), after adding
\(\langle\nabla\log u,A_\chi K(w,\cdot)\rangle\), is
\[
 d-3cvd-2(\partial_vc)v^2d+v\chi
 =\chi(2\chi-1).
\]
The mixed coefficient is \(1+d-2cv+v=2\chi\).
The remaining transverse terms are the last two terms of
\eqref{eq:fl-alg-exact-drift}; differentiating \(K\) produces its
first line. This proves the formula.

In particular every axial Hessian term in this divergence contains a
factor \(\chi\). Since \(H=S-K\) and
\(\tr P=F-\chi q_1\), the second line of
\eqref{eq:fl-alg-exact-drift} has absolute value at most
\[
 C|K|\bigl(|P^{\mathrm{tf}}|+|M|+|F|+\chi|q_1|\bigr)+C|K|^2
 \le\eta\mathcal T+C_\eta|K|^2.
\]
Here \(\chi|q_1|\le\sqrt{\chi}|q_1|\) and
\eqref{eq:fl-alg-floor-comparison} were used.
The first line is bounded by \(Ca|\nabla K|\).
Equation \eqref{eq:fl-alg-gradient-weight} proves
\eqref{eq:fl-alg-drift-derivative}, without an \(N\)-dependent constant.

\emph{The trace derivative in each stage.}
At fixed \((x,w,p)\), the height derivatives in
\eqref{eq:fl-alg-stage-right-sides} are respectively
\[
 1,\qquad 1+\theta\,\partial_hD_0,\qquad
 1+(1-\zeta)\partial_hD_0(x,w,h+\zeta b(x)),
\]
so each is at least one.
At \(\dd t=0\), the chain rule and \(h=\tau f\) give
\begin{equation}
 w(F)
 =\partial_h\mathfrak F\,\tau a|\dd f|
       +(\nabla_x\mathfrak F)(w)
       +D_w\mathfrak F(\nabla_w w),
 \qquad
 \nabla_w w=a(dH_{ee}e+H_{e\perp}).
 \label{eq:fl-alg-trace-chain}
\end{equation}
There is no \(\partial_p\mathfrak F\) term because \(\nabla p=0\).

The compact corrections have uniformly bounded first derivatives.
For the only noncompact correction, use
\[
 \tr_{A_\chi}K=\tr K-cK(w,w),\qquad
 D_wc[\xi]=2(\partial_vc)\langle w,\xi\rangle.
\]
The bounds
\[
 |c|\le C(1+N),\qquad
 |\partial_vc|\le C(1+N)^2
\]
give, on the stated height interval,
\[
 |\nabla_x\mathfrak F|
      \le C(1+N)^2\bigl(\mathbf 1_E+|\nabla K|\bigr),
 \qquad
 |D_w\mathfrak F|
      \le C(1+N)^2\bigl(\mathbf 1_E+|K|\bigr).
\]
The base derivative here also includes the bounded derivative of
\(b(x)\) inside \(D_0\).
The second term of \eqref{eq:fl-alg-trace-chain} is consequently
bounded below by
\(-C(1+N)^2(\rho_0+v\rho_1)\).
For its third term, the comparison
\eqref{eq:fl-alg-floor-comparison} and
\(d^2/\chi\le(N+4)/4\) yield
\[
 \begin{split}
 &C(1+N)^2(\mathbf 1_E+|K|)\,
       a\bigl(|H_{e\perp}|+d|H_{ee}|\bigr)\\
 &\hspace{12mm}\le
     \eta\mathcal T+
       C_\eta(1+N)^5(\mathbf 1_E+|K|^2).
 \end{split}
\]
For example, replace \(H\) by \(S-K\), use weighted Young Inequalities
against \(|M|^2+\chi q_1^2\), and absorb the remaining \(K\)-terms
into the same right side. Thus one may take \(m=5\), after enlarging
the constant. The weight observations at the start of the proof and
\(\partial_h\mathfrak F\ge1\) establish
\eqref{eq:fl-alg-trace-derivative}.

Finally, outside \(E\), all three right sides have the form
\[
 F=h+\zeta\tr_{A_\chi}K
\]
with \(0\le\zeta\le1\), where \(\zeta=0\) in the first two stages.
The eigenvalues of \(A_\chi\) lie in \((0,1]\), so
\(|F|\le |h|+C|K|\le Cr^{-\gamma}\) if the stated end height bound
holds. On \(E\), \(F\) is uniformly bounded by the compact hypotheses.
This proves \eqref{eq:fl-alg-weighted-remainder}.
\end{proof}

\begin{remark}
\label{rem:fl-alg-interfaces}
The preceding propositions and lemmas are identities and estimates
for smooth fields. Their application to the continuation uses the
right sides and compact corrections of
Proposition~\ref{prop:fl-homotopies}, the height bounds of
Lemma~\ref{lem:fl-height}, and the actual trace equation.
In particular no existence conclusion for the coupled system is
part of the scalar-curvature calculation.
\end{remark}
